\documentclass[a4paper,11pt]{article}
\usepackage[utf8]{inputenc}
\usepackage[T1]{fontenc}
\usepackage[ngerman]{babel}
\usepackage{geometry}
\geometry{left=2.5cm, right=2.5cm, top=2.5cm, bottom=2.5cm}
\usepackage{longtable}
\usepackage{booktabs}
\usepackage{array}
\usepackage{xcolor}
\usepackage{colortbl}
\usepackage{hyperref}
\usepackage{fancyhdr}
\usepackage{tabularx}
\usepackage{enumitem}

\definecolor{headerblue}{HTML}{1565C0}
\definecolor{tableheader}{HTML}{E3F2FD}
\definecolor{riskhigh}{HTML}{FFCDD2}
\definecolor{riskmedium}{HTML}{FFF9C4}
\definecolor{risklow}{HTML}{C8E6C9}

\hypersetup{colorlinks=true,linkcolor=headerblue,urlcolor=headerblue}

\pagestyle{fancy}
\fancyhf{}
\fancyhead[L]{\textbf{INSPECTAIR}}
\fancyhead[R]{Risk \& Opportunities v2.0}
\fancyfoot[C]{\thepage}
\renewcommand{\headrulewidth}{0.4pt}

\setlength{\parindent}{0pt}
\setlength{\parskip}{6pt}

\begin{document}

\begin{center}
{\LARGE\textbf{INSPECTAIR}}\\[0.3cm]
{\Large Risk \& Opportunities}\\[0.3cm]
{\large Risikobewertung und Chancenanalyse}\\[0.5cm]
\begin{tabular}{ll}
\textbf{Projekt:} & InspectAir -- Luftqualitätsmessgerät \\
\textbf{Version:} & 2.0 \\
\textbf{Datum:} & Februar 2026 \\
\textbf{Autoren:} & Maximilian Liebl, Cedric Stroh, Jannis Messer \\
\textbf{Modul:} & Mikrocomputertechnik 1, DHBW Ravensburg \\
\textbf{Dozent:} & Thomas Stingl (Airbus) \\
\end{tabular}
\end{center}

\vspace{0.5cm}
\noindent\textit{Dieses Dokument wurde unter Verwendung von KI-Tools (Claude, Anthropic) zur Überprüfung erstellt.}
\vspace{0.3cm}

\tableofcontents
\newpage

% ============================================================
\section{Risikoanalyse}

\subsection{R1 -- Stromspitzen durch 5V-Sensoren}

\begin{tabularx}{\textwidth}{lX}
\toprule
\rowcolor{riskhigh}
\textbf{Risikostufe} & \textbf{HOCH} \\
\midrule
Beschreibung & Der MH-Z19C und LD2410C erzeugen beim Einschalten und während der Messung kurzzeitige Stromspitzen bis zu 150\,mA. Diese können die 3{,}3V-Versorgung destabilisieren und zu Systemabstürzen oder verfälschten Messwerten führen. \\
Auswirkung & Systemabsturz, Fehlmessungen, unzuverlässiger Betrieb \\
Maßnahme & \textbf{220\,\textmu F Elektrolyt-Kondensatoren} an der VCC-Leitung beider 5V-Sensoren. Getestet und validiert (TC-08). \\
Status & \textbf{Mitigiert} -- Stützkondensatoren verbaut und getestet \\
\bottomrule
\end{tabularx}

\subsection{R2 -- 5V-Signale zerstören ESP32-GPIO}

\begin{tabularx}{\textwidth}{lX}
\toprule
\rowcolor{riskhigh}
\textbf{Risikostufe} & \textbf{HOCH} \\
\midrule
Beschreibung & Der LD2410C-Radar sendet UART-TX-Signale mit 5V-Pegel. Die ESP32-S3-GPIOs sind nur 3{,}3V-tolerant. Direktes Anschließen kann den GPIO-Pin oder den gesamten Mikrocontroller irreversibel beschädigen. \\
Auswirkung & Hardware-Zerstörung (ESP32-GPIO oder gesamter Chip) \\
Maßnahme & \textbf{TXS0108E Level Shifter} zwischen LD2410C-TX und ESP32-RX (GPIO7). Bidirektionaler Pegelwandler 3{,}3V $\leftrightarrow$ 5V. \\
Status & \textbf{Mitigiert} -- Level Shifter verbaut, PFLICHT-Bauteil \\
\bottomrule
\end{tabularx}

\subsection{R3 -- Zeitplan-Überschreitung}

\begin{tabularx}{\textwidth}{lX}
\toprule
\rowcolor{riskmedium}
\textbf{Risikostufe} & \textbf{MITTEL} \\
\midrule
Beschreibung & Komplexität der LVGL-9-Migration und Integration von 5 Sensoren plus 5 UI-Screens kann zu Verzögerungen führen. Abgabetermin: 13.03.2026. \\
Auswirkung & Verspätete Abgabe, unvollständige Dokumentation \\
Maßnahme & Wöchentliche Meilenstein-Reviews, Priorisierung der Kernfunktionen, parallele Dokumentation. \\
Status & \textbf{Aktiv überwacht} \\
\bottomrule
\end{tabularx}

\subsection{R4 -- Sensor-Ausfall}

\begin{tabularx}{\textwidth}{lX}
\toprule
\rowcolor{riskmedium}
\textbf{Risikostufe} & \textbf{MITTEL} \\
\midrule
Beschreibung & Einzelne Sensoren können defekt sein oder ausfallen (z.\,B. durch fehlerhafte Verkabelung, ESD-Schäden). \\
Auswirkung & Fehlende Messdaten für einzelne Parameter \\
Maßnahme & Ersatzsensoren beschafft, fehlerhafte Sensoren werden auf dem Display angezeigt (valid-Flag-Prüfung in jedem Lesezyklus), Watchdog-Timer verhindert Hänger. \\
Status & \textbf{Mitigiert} -- Fehlerbehandlung implementiert \\
\bottomrule
\end{tabularx}

\subsection{R5 -- I2C-Buskonflikt}

\begin{tabularx}{\textwidth}{lX}
\toprule
\rowcolor{risklow}
\textbf{Risikostufe} & \textbf{NIEDRIG} \\
\midrule
Beschreibung & AHT20 (Addr: 0x38) und SGP40 (Addr: 0x59) teilen sich den I2C-Bus. Gleichzeitiger Zugriff könnte Konflikte verursachen. \\
Auswirkung & Sporadische Lesefehler \\
Maßnahme & Sequentielle Abfrage in \texttt{aht\_sgp\_read()}, 4{,}7\,k$\Omega$ Pull-up-Widerstände, verschiedene I2C-Adressen. \\
Status & \textbf{Mitigiert} -- Kein Konflikt bei Tests festgestellt \\
\bottomrule
\end{tabularx}

\subsection{R6 -- Display-Lesbarkeit}

\begin{tabularx}{\textwidth}{lX}
\toprule
\rowcolor{risklow}
\textbf{Risikostufe} & \textbf{NIEDRIG} \\
\midrule
Beschreibung & 4-Zoll-Display könnte bei starkem Umgebungslicht schlecht ablesbar sein. \\
Auswirkung & Eingeschränkte Nutzbarkeit bei direkter Sonneneinstrahlung \\
Maßnahme & PWM-Backlight auf 100\,\% (255) im Active-Modus, vierstufige Farbcodierung mit hohem Kontrast, dunkler Hintergrund (Schwarz). \\
Status & \textbf{Akzeptiert} -- Innenraumgerät, keine direkte Sonneneinstrahlung vorgesehen \\
\bottomrule
\end{tabularx}

% ============================================================
\section{Chancenanalyse (Opportunities)}

\begin{tabularx}{\textwidth}{llX}
\toprule
\rowcolor{tableheader}
\textbf{ID} & \textbf{Chance} & \textbf{Beschreibung} \\
\midrule
O1 & WiFi-Datenlogging & Bereits implementiert: NTP-Zeitsynchronisation und Web-Fernsteuerung via HTTP-Server. Erweiterbar für Cloud-Anbindung (z.\,B. MQTT, InfluxDB). \\
O2 & Akkubetrieb & Light-Sleep-Modus ($\sim$0{,}8\,mA) bereits implementiert. Externe LiPo-Batterie mit Lademodul könnte portablen Einsatz ermöglichen. \\
O3 & UI-Themes & 5 verschiedene Screens bereits implementiert (Tree, Overview, Detail, Analog, Bubble). Weitere Themes über \texttt{ui\_manager} einfach ergänzbar. \\
O4 & Alarm-Funktionen & Grenzwert-Überschreitungen werden farblich signalisiert. Akustischer Alarm (Buzzer) oder Push-Benachrichtigungen wären einfach ergänzbar. \\
O5 & SD-Karten-Logging & Langzeit-Datenaufzeichnung auf SD-Karte möglich. SPI-Bus bereits für Display verwendet, zweiter SPI-Bus nutzbar. \\
\bottomrule
\end{tabularx}

% ============================================================
\section{Kostenschätzung}

\begin{tabularx}{\textwidth}{Xlr}
\toprule
\rowcolor{tableheader}
\textbf{Komponente} & \textbf{Anzahl} & \textbf{Preis (ca.)} \\
\midrule
ESP32-S3 N8R8 DevKit (Waveshare) & 1 & 13{,}00\,\euro \\
4" TFT Display ST7796S SPI & 1 & 15{,}00\,\euro \\
MH-Z19C CO\textsubscript{2}-Sensor (NDIR) & 1 & 22{,}00\,\euro \\
PMS5003 Feinstaub-Sensor & 1 & 18{,}00\,\euro \\
SGP40 VOC-Sensor & 1 & 8{,}00\,\euro \\
AHT20 Temperatur-/Feuchte-Sensor & 1 & 3{,}00\,\euro \\
LD2410C 24\,GHz Radar & 1 & 7{,}00\,\euro \\
\midrule
TXS0108E Level Shifter & 1 & 2{,}00\,\euro \\
220\,\textmu F Elektrolyt-Kondensatoren & 2 & 0{,}50\,\euro \\
BS170 MOSFET (Backlight) & 1 & 0{,}30\,\euro \\
4{,}7\,k$\Omega$ Widerstände (I2C Pull-ups) & 2 & 0{,}10\,\euro \\
10\,k$\Omega$ Widerstand (BS170 Gate) & 1 & 0{,}05\,\euro \\
PETG-Filament (3D-Druck) & -- & 5{,}00\,\euro \\
Breadboard + Kabel & -- & 5{,}00\,\euro \\
\midrule
\textbf{Gesamtkosten} & & \textbf{$\sim$99{,}00\,\euro} \\
\bottomrule
\end{tabularx}

% ============================================================
\section{Änderungshistorie}

\begin{tabularx}{\textwidth}{l l l X}
\toprule
\rowcolor{tableheader}
\textbf{Version} & \textbf{Datum} & \textbf{Autor} & \textbf{Änderung} \\
\midrule
1.0 & Januar 2026 & Team & Erstversion \\
2.0 & Februar 2026 & Team & Holzgehäuse $\to$ PETG 3D-Druck, Risiko-Status aktualisiert, Opportunities mit implementiertem Stand ergänzt, Kostenschätzung aktualisiert. \\
\bottomrule
\end{tabularx}

\end{document}
